\section{Discussion}
\label{sec:discussion}

\begin{table*}[t]
\centering
\caption{All registered predictions and their outcomes (values are seed 0/1/2 unless noted). Decision rules were fixed
before each run; the text reports the effect sizes.}
\label{tab:prereg}
\setlength{\tabcolsep}{3pt}
\footnotesize
\begin{tabular}{p{4.5cm}p{3.5cm}p{5.9cm}p{3.0cm}}
\toprule
Prediction & Decision rule & Observed & Outcome \\
\midrule
Name reaches the answer only in blocks 9--11 (full FT) & patching, $\Fl$ by layer & $\Fl(9) =$ \patchFLastLNineFullFT & built early, used late \\
Swap control (c-ii) keeps the shortcut (full FT) & kept $\ge 0.80$ & \swapCiiKeptLateFullFT & holds \\
Late blocks retrain better at 100 steps & late $-$ early $\ge 0.20$ & \budgetLateMinusEarlyFullFT & fails (holds in seed 0) \\
Swap and retrain maps agree & Spearman $\ge 0.6$ & FT \spearmanSwapRetrainFullFT; r4 \spearmanSwapRetrainLoraRFour; r64 \spearmanSwapRetrainLoraRSixtyFour & FT, r4 hold; r64 fails \\
Counter-examples do the fixing (3c) & $\ge 3$ blocks pass with names, fail without & FT \ccCountedBlocksFullFT, r4 \ccCountedBlocksLoraRFour, r64 \ccCountedBlocksLoraRSixtyFour{} blocks & holds for all three \\
Read-out at layer 9 removes the shortcut & removal criterion & \pSixRemovedLNineFullFT & holds \\
Random vector (same $\alpha$) removes $<$ half as much & ratio $< 0.5$ & \pSixRemovedRandomFullFT & fails (holds in seed 2) \\
Equal dose: layer 9 best & $l^*=9$ removes most & \pSixBEqLNineRemovedFullFT & holds \\
Random vector (own $\alpha$) removes $<$ half as much & ratio $< 0.5$ & \pSixBRandomRemovedFullFT & holds \\
Cue-free $\alpha$ still removes the shortcut & removal criterion & \pSixBCueFreeRemovedFullFT & holds \\
LoRA swap control (c-ii) keeps the shortcut & kept $\ge 0.80$ & r4 \swapCiiKeptLateLoraRFour; r64 \swapCiiKeptLateLoraRSixtyFour & fails \\
LoRA (c-ii) failure: init or training path (L5) & 3 replicate pairs per cell & A \lFiveKeptCellA; B \lFiveKeptCellB; C \lFiveKeptCellC & single-pair event \\
LoRA on blocks 0--5: arrival layer (L4) & $\ge 7$: GPT-2's structure & \lFourArrivalEarly & GPT-2's structure \\
Partly reliable name matches rank 4 (L6) & a level with gap 0.20--0.45 & best \lSixGapPNinetysZero{} (seed 0) & no level; stopped \\
Word-choice cue is learned (Phase 5) & gap $\ge 0.20$ & \gapWordcueFivesZero{} (5\%), \gapWordcueTensZero{} (10\%), seed 0 & not learned; stopped \\
Memorized labels: no single block removes them & no block passes & block \wlPassingBlockNoisy{} passes & fails \\
Pythia: lens reading (mean over seeds) & takeover layer & \pyTakeoverMeanPythia & decided early \\
\bottomrule
\end{tabular}
\end{table*}

\textbf{Find tells us where the answer is read out.} Three tools agree that in GPT-2 the name's information reaches the
answer only in blocks 9--11, while the true sentiment is still readable at layer~9; swaps show that these late blocks
need earlier computation. Placing LoRA adapters only on early blocks still yields a late arrival, so this is largely a
property of the pretrained network. In Pythia, lens and swaps agree on an early location. The tools are consistent
within a model, which is a precondition for using them to guide a fix.

\textbf{Counter-examples, not location, decide retraining.} Retraining any block removes the shortcut, but only when the
repair data shows the names with both labels. Every block writes into the shared residual stream, so any block can learn
to cancel the name's contribution once the data shows that it is uninformative. This explains, for a learned shortcut,
the mismatch between localization and editing that Hase et al.~\cite{hase2023localization} reported for facts. Two
further observations point the same way (both post-hoc): a few contradicting examples during training already weaken
the name shortcut sharply (Phase L6), and the word-choice cue, whose words appear with nearly balanced labels in the
untouched reviews, was not learned at all. Together they suggest that whether a shortcut forms, and whether it can be
trained away, depends on counter-examples in the data. Memorized wrong labels behave differently
(Sec.~\ref{sec:wronglabel}): with corrected labels only one early block fixes them, so for that habit the location does
matter.

\textbf{A training-free fix needs the location.} Reading the answer out at layer~9 removes the shortcut without training
data, works best at the layer the tools identify, is not reproduced by a matched random vector, and does not need to
know the cue. It is an early exit~\cite{halawi2024overthinking}, so its value here is not as a new method but as a test
of the question: the location found by interpretability tools does not tell us where to retrain, but it does tell us
where to read the answer out.

\textbf{LoRA.} Rank 64 reproduces the full fine-tuning picture, while rank 4 learns a weaker shortcut that only some
blocks can fix. The failed swap control under LoRA is associated with the training path rather than the adapter
initialization. Whether rank or shortcut strength drives the rank-4 exception remains open.
